% GENERATED FILE. Edit canonical lessons, then run npm run generate:books.
% canonical-source-hash: fnv1a64:e09579e5ed043bcb
% canonical-lessons: ML-C429-ksemam, ML-C429-resan, ML-C429-dhanasahayam, ML-C429-anukulyam, ML-C429-vitaranam-ceyyuka

\chapter{Welfare}
\label{ch:ml-welfare}
\hlchaptermodality{429}

\begin{chapteropening}
\textbf{By the end of this chapter:} \emph{I can say welfare, a ration, financial aid, a benefit and to distribute.}

Complete the last lesson of chapter 429: I can say welfare, a ration, financial aid, a benefit and to distribute.
\end{chapteropening}

\section[\texorpdfstring{\ml{ക്ഷേമം} (kṣēmaṁ)}{ക്ഷേമം (kṣēmaṁ)}]{\ml{ക്ഷേമം} (kṣēmaṁ) — welfare; well-being}

\label{lesson:ML-C429-ksemam}



\begin{quote}
Before the new one: say the Malayalam for a budget, then the Malayalam for statistics.
\end{quote}

\subsection*{You'll want to know: \ml{ക്ഷേമം}}
\textbf{\ml{ക്ഷേമം}} — \emph{kṣēmaṁ} — \textquotedblleft{}welfare; well-being\textquotedblright{}.

\begin{quote}
\ml{തൊഴിലാളികളുടെ} \ml{ക്ഷേമത്തിന്} \ml{പുതിയ} \ml{പദ്ധതി} \ml{വന്നു}.
\end{quote}

\begin{quote}
\emph{toḻilāḷikaḷuṭe kṣēmattinŭ putiya paddhati vannu.}
\end{quote}

\begin{quote}
A new scheme has come for the workers' welfare.
\end{quote}

As a greeting in letters: \ml{ക്ഷേമം} \ml{തന്നെയല്ലേ}, all well, I hope.

\begin{cousinweb}
\textbf{\ml{ക്ഷേമം}} (\emph{kṣēmaṁ}) is borrowed from Sanskrit \emph{kṣema}, \textquotedblleft{}giving rest or ease; safety, well-being\textquotedblright{} (Monier-Williams).
\end{cousinweb}

\begin{grammarlens}[title={What you've built}]
One more word for giving reasons and opinions.
\end{grammarlens}

\subsection*{Your turn}
\begin{itemize}
\raggedright
  \item \emph{Say it:} \emph{kṣēmaṁ}
  \item \emph{Say it:} \emph{kṣēmaṁ}, once more
  \item \emph{Say it:} \emph{toḻilāḷikaḷuṭe kṣēmattinŭ putiya paddhati vannu.} — A new scheme has come for the workers' welfare.
  \item \emph{Recall:} say the Malayalam for a budget, then the Malayalam for statistics, then say \emph{kṣēmaṁ} again
\end{itemize}

\begin{tcolorbox}[breakable,colback=teal!4,colframe=teal!35!black,title={Before you move on}]
What does \ml{ക്ഷേമം} mean? (\textquotedblleft{}Welfare; well-being\textquotedblright{}.) Say it once more.
\end{tcolorbox}
\section[\texorpdfstring{\ml{റേഷൻ} (ṟēṣan)}{റേഷൻ (ṟēṣan)}]{\ml{റേഷൻ} (ṟēṣan) — a ration}

\label{lesson:ML-C429-resan}



\begin{quote}
Before the new one: say the Malayalam for statistics, then the Malayalam for welfare.
\end{quote}

\subsection*{You'll want to know: \ml{റേഷൻ}}
\textbf{\ml{റേഷൻ}} — \emph{ṟēṣan} — \textquotedblleft{}a ration\textquotedblright{}.

\begin{quote}
\ml{റേഷൻ} \ml{കടയിൽ} \ml{ഇന്ന്} \ml{അരി} \ml{വന്നു}.
\end{quote}

\begin{quote}
\emph{ṟēṣan kaṭayil innŭ ari vannu.}
\end{quote}

\begin{quote}
Rice came to the ration shop today.
\end{quote}

The ration shop, \ml{റേഷൻ} \ml{കട}, sells rice and sugar at fixed prices.

\begin{cousinweb}
\textbf{\ml{റേഷൻ}} (\emph{ṟēṣan}) was borrowed from English \emph{ration}.
\end{cousinweb}

\begin{grammarlens}[title={What you've built}]
One more word for giving reasons and opinions.
\end{grammarlens}

\subsection*{Your turn}
\begin{itemize}
\raggedright
  \item \emph{Say it:} \emph{ṟēṣan}
  \item \emph{Say it:} \emph{ṟēṣan}, once more
  \item \emph{Say it:} \emph{ṟēṣan kaṭayil innŭ ari vannu.} — Rice came to the ration shop today.
  \item \emph{Recall:} say the Malayalam for statistics, then the Malayalam for welfare, then say \emph{ṟēṣan} again
\end{itemize}

\begin{tcolorbox}[breakable,colback=teal!4,colframe=teal!35!black,title={Before you move on}]
What does \ml{റേഷൻ} mean? (\textquotedblleft{}A ration\textquotedblright{}.) Say it once more.
\end{tcolorbox}
\section[\texorpdfstring{\ml{ധനസഹായം} (dhanasahāyaṁ)}{ധനസഹായം (dhanasahāyaṁ)}]{\ml{ധനസഹായം} (dhanasahāyaṁ) — financial aid}

\label{lesson:ML-C429-dhanasahayam}



\begin{quote}
Before the new one: say the Malayalam for welfare, then the Malayalam for a ration.
\end{quote}

\subsection*{You'll want to know: \ml{ധനസഹായം}}
\textbf{\ml{ധനസഹായം}} — \emph{dhanasahāyaṁ} — \textquotedblleft{}financial aid\textquotedblright{}.

\begin{quote}
\ml{വീട്} \ml{നഷ്ടപ്പെട്ടവർക്ക്} \ml{സർക്കാർ} \ml{ധനസഹായം} \ml{നൽകി}.
\end{quote}

\begin{quote}
\emph{vīṭŭ naṣṭappeṭṭavarkkŭ sarkkār dhanasahāyaṁ nalki.}
\end{quote}

\begin{quote}
The government gave financial aid to those who lost their homes.
\end{quote}

\ml{സഹായം} alone is help.

\begin{cousinweb}
\textbf{\ml{ധനസഹായം}} (\emph{dhanasahāyaṁ}) is built from Sanskrit \emph{dhana}, \textquotedblleft{}property, wealth\textquotedblright{} (CDIAL 6717), and \emph{sahāya}, \textquotedblleft{}a companion, a helper\textquotedblright{} (CDIAL 13309): help in money.
\end{cousinweb}

\begin{grammarlens}[title={What you've built}]
One more word for giving reasons and opinions.
\end{grammarlens}

\subsection*{Your turn}
\begin{itemize}
\raggedright
  \item \emph{Say it:} \emph{dhanasahāyaṁ}
  \item \emph{Say it:} \emph{dhanasahāyaṁ}, once more
  \item \emph{Say it:} \emph{vīṭŭ naṣṭappeṭṭavarkkŭ sarkkār dhanasahāyaṁ nalki.} — The government gave financial aid to those who lost their homes.
  \item \emph{Recall:} say the Malayalam for welfare, then the Malayalam for a ration, then say \emph{dhanasahāyaṁ} again
\end{itemize}

\begin{tcolorbox}[breakable,colback=teal!4,colframe=teal!35!black,title={Before you move on}]
What does \ml{ധനസഹായം} mean? (\textquotedblleft{}Financial aid\textquotedblright{}.) Say it once more.
\end{tcolorbox}
\section[\texorpdfstring{\ml{ആനുകൂല്യം} (ānukūlyaṁ)}{ആനുകൂല്യം (ānukūlyaṁ)}]{\ml{ആനുകൂല്യം} (ānukūlyaṁ) — a benefit; a concession}

\label{lesson:ML-C429-anukulyam}



\begin{quote}
Before the new one: say the Malayalam for a ration, then the Malayalam for financial aid.
\end{quote}

\subsection*{You'll want to know: \ml{ആനുകൂല്യം}}
\textbf{\ml{ആനുകൂല്യം}} — \emph{ānukūlyaṁ} — \textquotedblleft{}a benefit; a concession\textquotedblright{}.

\begin{quote}
\ml{പ്രായമായവർക്ക്} \ml{ബസിൽ} \ml{ചില} \ml{ആനുകൂല്യങ്ങൾ} \ml{ഉണ്ട്}.
\end{quote}

\begin{quote}
\emph{prāyamāyavarkkŭ basil cila ānukūlyaṅṅaḷ uṇṭŭ.}
\end{quote}

\begin{quote}
Elderly people get some concessions on the bus.
\end{quote}

Plural: \ml{ആനുകൂല്യങ്ങൾ}, benefits.

\begin{cousinweb}
\textbf{\ml{ആനുകൂല്യം}} (\emph{ānukūlyaṁ}) is borrowed from Sanskrit \emph{ānukūlya}, \textquotedblleft{}conformity, suitableness, favour\textquotedblright{} (Monier-Williams), made from \emph{anukūla}, \textquotedblleft{}favourable\textquotedblright{}.
\end{cousinweb}

\begin{grammarlens}[title={What you've built}]
One more word for giving reasons and opinions.
\end{grammarlens}

\subsection*{Your turn}
\begin{itemize}
\raggedright
  \item \emph{Say it:} \emph{ānukūlyaṁ}
  \item \emph{Say it:} \emph{ānukūlyaṁ}, once more
  \item \emph{Say it:} \emph{prāyamāyavarkkŭ basil cila ānukūlyaṅṅaḷ uṇṭŭ.} — Elderly people get some concessions on the bus.
  \item \emph{Recall:} say the Malayalam for a ration, then the Malayalam for financial aid, then say \emph{ānukūlyaṁ} again
\end{itemize}

\begin{tcolorbox}[breakable,colback=teal!4,colframe=teal!35!black,title={Before you move on}]
What does \ml{ആനുകൂല്യം} mean? (\textquotedblleft{}A benefit; a concession\textquotedblright{}.) Say it once more.
\end{tcolorbox}
\section[\texorpdfstring{\ml{വിതരണം} \ml{ചെയ്യുക} (vitaraṇaṁ ceyyuka)}{വിതരണം ചെയ്യുക (vitaraṇaṁ ceyyuka)}]{\ml{വിതരണം} \ml{ചെയ്യുക} (vitaraṇaṁ ceyyuka) — to distribute}

\label{lesson:ML-C429-vitaranam-ceyyuka}



\begin{quote}
Before the new one: say the Malayalam for financial aid, then the Malayalam for a benefit.
\end{quote}

\subsection*{You'll want to know: \ml{വിതരണം} \ml{ചെയ്യുക}}
\textbf{\ml{വിതരണം} \ml{ചെയ്യുക}} — \emph{vitaraṇaṁ ceyyuka} — \textquotedblleft{}to distribute\textquotedblright{}.

\begin{quote}
\ml{അധ്യാപകർ} \ml{കുട്ടികൾക്ക്} \ml{പുസ്തകങ്ങൾ} \ml{വിതരണം} \ml{ചെയ്തു}.
\end{quote}

\begin{quote}
\emph{adhyāpakar kuṭṭikaḷkkŭ pustakaṅṅaḷ vitaraṇaṁ ceytu.}
\end{quote}

\begin{quote}
The teachers distributed books to the children.
\end{quote}

Past: \ml{വിതരണം} \ml{ചെയ്തു}, \emph{vitaraṇaṁ ceytu}.

The noun \ml{വിതരണം} is distribution.

\begin{cousinweb}
\textbf{\ml{വിതരണം} \ml{ചെയ്യുക}} (\emph{vitaraṇaṁ ceyyuka}) is a phrase verb built on Sanskrit \emph{vitaraṇa}, \textquotedblleft{}giving away, granting\textquotedblright{} (Monier-Williams), with \ml{ചെയ്യുക}, \textquotedblleft{}to do\textquotedblright{}.
\end{cousinweb}

\begin{grammarlens}[title={What you've built}]
One more word for giving reasons and opinions.
\end{grammarlens}

\subsection*{Your turn}
\begin{itemize}
\raggedright
  \item \emph{Say it:} \emph{vitaraṇaṁ ceyyuka}
  \item \emph{Say it:} \emph{vitaraṇaṁ ceyyuka}, once more
  \item \emph{Say it:} \emph{adhyāpakar kuṭṭikaḷkkŭ pustakaṅṅaḷ vitaraṇaṁ ceytu.} — The teachers distributed books to the children.
  \item \emph{Recall:} say the Malayalam for financial aid, then the Malayalam for a benefit, then say \emph{vitaraṇaṁ ceyyuka} again
\end{itemize}

\begin{tcolorbox}[breakable,colback=teal!4,colframe=teal!35!black,title={Before you move on}]
What does \ml{വിതരണം} \ml{ചെയ്യുക} mean? (\textquotedblleft{}To distribute\textquotedblright{}.) Say it once more.
\end{tcolorbox}
